%!TEX program = xelatex
\documentclass[lang=en,11pt,a4paper,citestyle=authoryear]{elegantpaper}

\title{HW3 - STAT 709}
\author{Wells Guan}

\usepackage{array,url,stix}
\usepackage{tikz}
\usepackage{tikz-cd}
\usepackage{enumitem}
\usepackage{amsmath}
\usepackage{amssymb}
\usepackage{listings}
\usepackage{subfigure}

\lstset{
    frame=single,
    columns=fullflexible,
    keepspaces=true,
    showstringspaces=false,
    breaklines=true,
    breakatwhitespace=false,
    basicstyle=\ttfamily\footnotesize,
    numbers=left,
    numberstyle=\tiny,
    numbersep=6pt
}
\newcommand{\Z}{\mathbb{Z}}
\newcommand{\R}{\mathbb{R}}
\newcommand{\N}{\mathbb{N}}
\newcommand{\B}{\mathcal{B}}
\newcommand{\E}{\mathbb{E}}
\renewcommand{\P}{\mathbb{P}}
\newcommand{\norm}[1]{\left\lVert#1\right\rVert}

\begin{document}
\maketitle

% Textbook statements are not included in HW3_assignment.pdf.
% Insert their exact wording from the indicated editions when available.
\subsection*{Problem 1. Textbook exercises}
Jun Shao's book, Chapter 1: Exercises 50 and 53 (required); Exercise 91 (optional).

\subsubsection*{Ex. 1.60}
Let $X$ be a random variable with $\E X^2 < \infty$ and let $Y = |X|$. Suppose that $X$ has a Lebesgue p.d.f. symmetric about $0$. Show that $X$ and $Y$ are uncorrelated, but they are not independent.
\par\vspace{0.5em}
\textbf{Sol.}\par
Notice that
\[
I:= \int_0^{\infty} xf(x)dx = -\int_{-\infty}^0 x f(x)dx
\]
and hence $\E X = 0$ and $\E Y = 2I$. Assume $f$ to be the p.d.f. of $X$, then
\[
\begin{aligned}
    \text{Cov}(X,Y) &= \E XY \\
    &= \int_0^{\infty} x^2 f(x)dx -\int_{-\infty}^0 x^2f(x)dx = 0  
\end{aligned}
\]
since $f(x) = f(-x)$ for any $x\in\R$, which means $X$ and $Y$ are uncorrelated.\par
Notice that
\[
P(X \geq a, Y\geq a) = P(X\geq a), \quad P(Y\geq a) = 2P(X\geq a)  
\]
and if $X$ and $Y$ are independent, then we have
\[2P(X\geq a)^2 = P(X\geq a),\] which means that $P(X\geq a) = 1/2$ for any $a \in \R$. This is a contradiction to $EX^2 < \infty$, since
\[ \E X^2 \geq 2a^2P(X\geq a)\]
for any $a\in \R$.
\vspace{1em}

\subsubsection*{Ex. 1.82}

Let $(X_1,X_2)$ be $N_k(\mu,\Sigma)$ with a $k\times k$
positive definite
\[
\Sigma=
\begin{pmatrix}
    \Sigma_{11} & \Sigma_{12}\\
    \Sigma_{21} & \Sigma_{22}
\end{pmatrix},
\]
where $X_1$ is a random $l$-vector and $\Sigma_{11}$ is an
$l\times l$ matrix. Show that the conditional Lebesgue p.d.f.
of $X_2$ given $X_1=x_1$ is
\[
N_{k-l}\left(
    \mu_2+\Sigma_{21}\Sigma_{11}^{-1}(x_1-\mu_1),
    \Sigma_{22}-\Sigma_{21}\Sigma_{11}^{-1}\Sigma_{12}
\right),
\]
where $\mu_i=\E X_i$, $i=1,2$.
\par\vspace{0.5em}
\textbf{Sol.}\par
We would like to prove that
\[
\mu((\mathbf{x}_1,\mathbf{x}_2),A) := \int_A \dfrac{1}{(2\pi)^{(k-l)/2}\left(\det\left(\Sigma'\right)\right)^{1/2}}\exp\left(-\dfrac{1}{2}(\mathbf{x} - \mu'(\mathbf{x}_1))^T\Sigma'^{-1}(\mathbf{x} - \mu'(\mathbf{x}_1))\right) d\mathbf{x}
\]
for $\mathbf{x}_1\in\R^l, \mathbf{x}_2 \in \R^{k-l}$ and where
\[
\Sigma':=\Sigma_{22}-\Sigma_{21}\Sigma_{11}^{-1}\Sigma_{12},\quad \mu'(\mathbf{x}_1):= \mu_2+\Sigma_{21}\Sigma_{11}^{-1}(\mathbf{x}_1-\mu_1)
\]
is a regular conditional distribution from $\R^k\times \B^{k-l} \to [0,1]$. It remains to show that $\mu(\cdot,A)$ is a version of $P(X_2 \in A|X_1)$, and notice
\[
\begin{aligned}
    &\int_{\{X_1 \leq a\}} P(X_2 \in A|X_1)dP\\
    =& \int_{\mathbf{x}_1 \leq a}\int_{\mathbf{x}_2 \in A}\dfrac{1}{(2\pi)^{k/2}\left(\det \Sigma\right)^{1/2}}\exp\left(-\dfrac{1}{2}((\mathbf{x_1},\mathbf{x_2})-(\mu_1,\mu_2))^T\Sigma^{-1}((\mathbf{x_1},\mathbf{x_2})-(\mu_1,\mu_2))\right)d\mathbf{x_1}d\mathbf{x_2} \\
    =& \int_{\mathbf{x}_1\leq a}\int_A\dfrac{1}{(2\pi)^{(k-l)/2}\left(\det\left(\Sigma'\right)\right)^{1/2}}\exp\left(-\dfrac{1}{2}(\mathbf{x} - \mu'(\mathbf{x}_1))^T\Sigma'^{-1}(\mathbf{x} - \mu'(\mathbf{x}_1))\right) d\mathbf{x}\\
    &\dfrac{1}{(2\pi)^{l/2}\left(\det \Sigma_{11}\right)^{1/2}}\exp\left(-\dfrac{1}{2}(\mathbf{x_1}-\mu_1)^T\Sigma_{11}^{-1}(\mathbf{x_1}-\mu_1)\right)d\mathbf{x}_1\\
    =& \int\int_{\mathbf{x}_1\leq a}\int_A \dfrac{1}{(2\pi)^{(k-l)/2}\left(\det\left(\Sigma'\right)\right)^{1/2}}\exp\left(-\dfrac{1}{2}(\mathbf{x} - \mu'(\mathbf{x}_1))^T\Sigma'^{-1}(\mathbf{x} - \mu'(\mathbf{x}_1))\right) d\mathbf{x}\\
    &\dfrac{1}{(2\pi)^{k/2}\left(\det \Sigma\right)^{1/2}}\exp\left(-\dfrac{1}{2}((\mathbf{x_1},\mathbf{x_2})-(\mu_1,\mu_2))^T\Sigma^{-1}((\mathbf{x_1},\mathbf{x_2})-(\mu_1,\mu_2))\right)d\mathbf{x}_1d\mathbf{x}_2 \\
    &= \int_{\{X_1\leq a\}}\mu(\cdot,A)dP
\end{aligned}
\]
where
\[
\det \Sigma' \det\Sigma_{11} = \det \Sigma
\]
by Schur's lemma and that
\[
\Sigma^{-1}=
\begin{pmatrix}
    \Sigma_{11}^{-1} + \Sigma_{11}^{-1}\Sigma_{12}\Sigma'^{-1}\Sigma_{21}\Sigma_{11}^{-1} & -\Sigma_{11}^{-1}\Sigma_{12}\Sigma'^{-1}\\
    -\Sigma'^{-1}\Sigma_{21}\Sigma_{11}^{-1} & \Sigma'^{-1}
\end{pmatrix},
\]
for any $a \in \R$ and we are done.
\vspace{1em}

\subsubsection*{Ex. 1.91 (optional)}
% TODO: Insert the textbook problem statement.
\textit{[Problem statement to be added from the textbook.]}
\par\vspace{0.5em}
\textbf{Sol.}\par
% TODO: Write solution.
\vspace{1em}

\subsection*{Problem 2. Textbook exercises}
Durrett's \textit{Probability: Theory and Examples}, fifth edition: Exercises 2.1.4 and 2.1.6.

\subsubsection*{Ex. 2.1.4}
Let $\Omega = (0,1), \mathcal{F} =$ Borel sets, $P =$ Lebesgue measure. $X_n(\omega) =
\sin(2\pi n\omega), n = 1,2,...$ are uncorrelated but not independent.
\par\vspace{0.5em}
\textbf{Sol.}\par
Notice that
\[
\begin{aligned}
\text{Cov}(X_m,X_n) &= \int_0^1 \sin(2\pi m \omega)\sin(2\pi n \omega)d\omega \\ &= \dfrac{1}{2}\left(\int_0^1 \cos(2\pi(m-n)\omega)d\omega - \int_0^1 \cos(2\pi(m+n)\omega)d\omega\right)
\\ &= 0 \end{aligned}
\]
for $m\neq n$, and hence they are uncorrelated. To see they are not independent, we may notice that
\[
P( -1 \leq X_n \leq -1/2,  -1 \leq X_{2n} \leq -1/2) = 1/6 \neq 1/9 = P(-1 \leq X_n \leq -1/2)P(-1 \leq X_{2n} \leq -1/2)
\]
for $n\geq 1$ and we are done.
\vspace{1em}

\subsubsection*{Ex. 2.1.6}
Prove directly from the definition that if $X$ and $Y$ are independent and $f$ and $g$ are measurable functions then $f(X)$ and $g(Y)$ are
independent.
\par\vspace{0.5em}
\textbf{Sol.}\par
For any $A,B \in \B(\R)$, we know
\[
\begin{aligned}
P(f(X) \in A, g(Y)\in B) &= P(X\in f^{-1}(A), Y\in g^{-1}(B))\\ &= P(X\in f^{-1}(A))P(Y\in g^{-1}(B))\\ &= P(f(X) \in A)P(g(Y) \in B)
\end{aligned}
\]
and we are done.
\vspace{1em}

\newpage
\noindent\textit{The remaining exercises concern conditional expectation and are optional. Read the questions and try as many parts as you can.}

\subsection*{Problem 3. Conditioning on a transformation (optional)}
Let $X\in\R^n$ have a continuous density $p$, and let $f,g:\R^n\to\R^n$ be continuous, with $\E\norm{f(X)}_2<\infty$.
Let $A_1,\ldots,A_m$ be pairwise disjoint open sets whose union has a Lebesgue-null complement. On each $A_j$, assume that $g$ is continuously differentiable and one-to-one, with $\det Dg(x)\neq0$. Write $V_j=g(A_j)$ and let $h_j:V_j\to A_j$ be its inverse.

We want a formula for $\E[f(X)\mid g(X)=y]$. In general, $(f(X),g(X))$ does not have a joint density on $\R^{2n}$, so the elementary joint-density ratio formula need not apply directly.

For the following local calculation, fix $y\notin\bigcup_{j=1}^m\partial V_j$. Put $J(y)=\{j:y\in V_j\}$, and assume that $p(h_j(y))>0$ for at least one $j\in J(y)$. These conditions ensure that a sufficiently small cube around $y$ meets exactly the active inverse branches and that the limiting denominator is positive. For $\Delta>0$, let $E_\Delta=\prod_{i=1}^n(y_i-\Delta,y_i+\Delta)$.

\begin{enumerate}[label=(\arabic*),leftmargin=*]
\item For $j\in J(y)$ and sufficiently small $\Delta$, define $B_\Delta^{(j)}=h_j(E_\Delta)$. Prove that
\[
\operatorname{Vol}(B_\Delta^{(j)})=|\det Dh_j(y)|\,(2\Delta)^n+o(\Delta^n)\qquad(\Delta\downarrow0).
\]
\item Use part (1) to evaluate
\[
\lim_{\Delta\downarrow0}\frac{\E[f(X)\mathbf{1}\{g(X)\in E_\Delta\}]}{\P(g(X)\in E_\Delta)}.
\]
Derive a version of $\E[f(X)\mid g(X)=y]$ in terms of the inverse branches. Explain why this formula, summed over the branches defined at $y$, gives a conditional expectation for almost every $y$ under the law of $g(X)$.
\end{enumerate}
\vspace{0.5em}
\textbf{Sol.}\par
\begin{enumerate}[label=(\arabic*),leftmargin=*]
\item We know\[\begin{aligned}
\text{Vol}(B_{\Delta}^{(j)}) =& \int_{h_j(E_{\Delta})} d\mathbf{x} \\=& \int_{E_{\Delta}}|\det Dh_j|(\mathbf{y})d\mathbf{y} \\=& |\det Dh_j(y)|(2\Delta)^n + \int_{E_{\Delta}}\left(|\det Dh_j(\mathbf{y})| - |\det Dh_j(y)|\right)d\mathbf{y}
\end{aligned}
\]
where
\[
\begin{aligned}
&\lim\limits_{\Delta\to 0}\left|\int_{E_{\Delta}}\left(|\det Dh_j(\mathbf{y})| - |\det Dh_j(y)|\right)d\mathbf{y}\right|/(\Delta)^n \\ \leq& 4^n\lim_{\Delta \to 0}\max_{E_{\Delta}}||\det Dh_j(\mathbf{y})| - |\det Dh_j(y)|| = 0
\end{aligned}\]
and we are done since $h_j$ is continuous.
\item Notice that
\[
\begin{aligned}
\E[f(X)\mathbf{1}\{g(X)\in E_\Delta\}] &= \int_{\bigcup_{j\in J}h_j(E_{\Delta})} f(\mathbf{x})p(\mathbf{x})d\mathbf{x} \\
&=\sum\limits_{j\in J}\int_{E_{\Delta}} f(h_j(\mathbf{y}))p(h_j(\mathbf{y}))|\det Dh_j|(\mathbf{y})d\mathbf{y} \\
&= \sum\limits_{j\in J} f(h_j(y))p(h_j(y))|\det Dh_j(y)|(2\Delta)^n + o(\Delta^n)
\end{aligned}\]
with a similar proof in (1) and
\[
\P(g(X)\in E_\Delta) = \int_{\bigcup_{j\in J}h_j(E_{\Delta})} p(\mathbf{x})d\mathbf{x} = \sum\limits_{j\in J}p(h_j(y))|\det Dh_j(y)|\,(2\Delta)^n+o(\Delta^n)
\]
and hence
\[
\begin{aligned}
    \lim_{\Delta\downarrow0}\frac{\E[f(X)\mathbf{1}\{g(X)\in E_\Delta\}]}{\P(g(X)\in E_\Delta)} &= \lim_{\Delta\downarrow0}\dfrac{\sum\limits_{j\in J} f(h_j(y))p(h_j(y))|\det Dh_j(y)|(2\Delta)^n + o(\Delta^n)}{\sum\limits_{j\in J}p(h_j(y))|\det Dh_j(y)|\,(2\Delta)^n+o(\Delta^n)} \\
    &= \dfrac{\sum\limits_{j\in J} f(h_j(y))p(h_j(y))|\det Dh_j(y)|}{\sum\limits_{j\in J}p(h_j(y))|\det Dh_j(y)|}
\end{aligned}
\]
which is a version of $\E[f(X)\mid g(X)=y]$. To see this, it sufficents to prove that
\[
\E[f(X)|g(X)] = \dfrac{\sum\limits_{j\in J(g(X))} f(h_j(g(X)))p(h_j(g(X)))|\det Dh_j(g(X))|}{\sum\limits_{j\in J(g(X))}p(h_j(g(X)))|\det Dh_j(g(X))|}
\]
which can be shown by noticing
\[
\begin{aligned}
    \int_{g(X) \in A} f(X)dP &= \int_{g(\mathbf{x}) \in A} f(\mathbf{x})p(\mathbf{x})d\mathbf{x} \\
    & = \sum\limits_{k=1}^m\int_{\mathbf{y}\in A\cap V_k'} f(h_k(\mathbf{y}))p(h_k(\mathbf{y}))|\det Dh_k(\mathbf{y})|d\mathbf{y} 
    \\
    & = \int_{g(\mathbf{x})\in A} \dfrac{\sum\limits_{j\in J(g(\mathbf{x}))} f(h_j(g(\mathbf{x})))p(h_j(g(\mathbf{x})))|\det Dh_j(g(\mathbf{x}))|}{\sum\limits_{j\in J(g(\mathbf{x}))}p(h_j(g(\mathbf{x})))|\det Dh_j(g(\mathbf{x}))|} p(\mathbf{x})d\mathbf{x}
    \\&= \int_{g(X) \in A} \dfrac{\sum\limits_{j\in J(g(X))} f(h_j(g(X)))p(h_j(g(X)))|\det Dh_j(g(X))|}{\sum\limits_{j\in J(g(X))}p(h_j(g(X)))|\det Dh_j(g(X))|} dP
\end{aligned}
\]
\end{enumerate}

\newpage
\subsection*{Problem 4. Conditional independence (optional)}
Let $X,Y,Z$ be real-valued random variables. Assume that $X$ and $Y$ are conditionally independent given $Z$, in the sense that
\[
\P(B\mid X,Z)=\P(B\mid Z)\quad\text{a.s. for every }B\in\sigma(Y).
\]
The first two parts connect this definition with the familiar density formulation.

\begin{enumerate}[label=(\arabic*),leftmargin=*]
\item Let $(U,W)$ be random vectors with a joint density. For a Borel function $\varphi$ satisfying $\E|\varphi(U)|<\infty$, show that a version is
\[
\E[\varphi(U)\mid W=w]=\int\varphi(u)p(u\mid w)\,du
\]
for almost every $w$ under the law of $W$, where $p(u\mid w)=p_{U,W}(u,w)/p_W(w)$ when $p_W(w)>0$.

\item If $(X,Y,Z)$ has a joint density, show that
\[
f_{X,Y\mid Z}(x,y\mid z)=f_{X\mid Z}(x\mid z)f_{Y\mid Z}(y\mid z)
\]
for almost every $z$ under the law of $Z$ and Lebesgue-almost every $(x,y)$.

\item Without assuming a joint density, prove that, for every $A\in\sigma(X)$ and $B\in\sigma(Y)$,
\[
\P(A\cap B\mid Z)=\P(A\mid Z)\P(B\mid Z)\qquad\text{a.s.}
\]
\textit{Hint:} Use the tower property and the definition of conditional expectation.

\item Let $\mathbf{X}=(X_1,\ldots,X_p)^\top\sim N_p(0,\Sigma)$, with $\Sigma$ positive definite, and define the precision matrix $\Theta=\Sigma^{-1}$. Set $Y=X_1$ and $Z=(X_2,\ldots,X_p)^\top$. In the coordinate order $(Y,Z)$, write
\[
\Sigma=\begin{pmatrix}\sigma_{YY}&\sigma_{ZY}^\top\\\sigma_{ZY}&\Sigma_{ZZ}\end{pmatrix}.
\]
Recall that
\[
Y\mid Z=z\sim N\!\left(z^\top\Sigma_{ZZ}^{-1}\sigma_{ZY},\,
\sigma_{YY}-\sigma_{ZY}^\top\Sigma_{ZZ}^{-1}\sigma_{ZY}\right),
\qquad\theta_{ZY}=-\theta_{YY}\Sigma_{ZZ}^{-1}\sigma_{ZY}.
\]
In a Gaussian graphical model, distinct nodes are joined when the corresponding precision-matrix entry is \textit{nonzero}; a zero entry indicates a missing edge. Prove that $\theta_{ij}=0$ implies that $X_i$ and $X_j$ are conditionally independent given all the other coordinates.
\end{enumerate}
\vspace{0.5em}
\textbf{Sol.}\par
\begin{enumerate}[label=(\arabic*),leftmargin=*]
\item % TODO: Write solution.
\item % TODO: Write solution.
\item % TODO: Write solution.
\item % TODO: Write solution.
\end{enumerate}

\end{document}
